\subsection{A Rule for the Season}\label{sec:rule}

Watching, prayer, fasting and alms can be gathered into a rule: a short written
plan of what one will do each day and each week until Christmas. A rule turns good
intentions into a few definite acts, and it makes them easy to review. Below
are three measures of the same rule. Choose one, keep it whole, and do not add
to it in the middle of the season.

\textit{A beginning} suits someone who has not kept Advent before, or whose
duties leave little room. \textit{A steady rule} suits someone who already
prays each day and wants the season to shape the day. \textit{A full rule}
suits someone with more freedom and some experience of the spiritual life; it
is best taken up with the advice of a confessor or director, who knows the
person and the duties it must fit.

\medskip
{\small
\begin{longtable}{@{}>{\raggedright\arraybackslash}p{0.13\linewidth}>{\raggedright\arraybackslash}p{0.255\linewidth}>{\raggedright\arraybackslash}p{0.255\linewidth}>{\raggedright\arraybackslash}p{0.255\linewidth}@{}}
\toprule
& A beginning & A steady rule & A full rule\\
\midrule
\endfirsthead
\toprule
& A beginning & A steady rule & A full rule\\
\midrule
\endhead
Prayer &
The Sunday's Collect each morning; the Angelus. &
As the beginning, with the \latin{Alma Redemptoris Mater} at night and the O
antiphon from 17 December. &
As the steady rule, with Vespers of the Little Office in its Advent form, or
the hymn \latin{Creator alme siderum} at evening.\\
\addlinespace
Scripture &
The Sunday Epistle and Gospel, read on Saturday evening. &
The day's Isaiah and the Father beside it (Appendix~\ref{sec:daily}). &
The day's Isaiah read in the four steps of the ladder, and the novena of
meditations from 16 December.\\
\addlinespace
Mass &
Sunday Mass, and the Immaculate Conception on 8 December. &
As the beginning, with Mass on one of the Ember Days. &
Daily Mass where it can be had.\\
\addlinespace
Fasting &
Fridays kept as days of penance. &
As the beginning, with the Ember Days and the Vigil kept as fasts. &
As the steady rule, with the Wednesdays and Fridays of Advent kept with fast
and abstinence.\\
\addlinespace
Alms &
One gift to the poor, given before Christmas. &
The money the fast saves, given away each week. &
As the steady rule, with one work of mercy done in person each week.\\
\addlinespace
Examen and confession &
The examen at night; confession in the week of the Ember Days. &
As the beginning, with the examination by the weeks of Advent in the nights
before confession. &
As the steady rule, with a fixed time of silence each day.\\
\bottomrule
\end{longtable}
}

Write the rule down on the evening before the First Sunday, on one side of a
card, and keep it where you pray. Read it once a week, on Saturday evening, and
ask of the week just past what was kept, what was dropped, and why.

A rule fails in two ways. The first is too much: a plan drawn up in the fervour
of the First Sunday that collapses by the second week. Choose less than you
think you can do. The second is drift: the rule is kept for a few days and then
quietly forgotten. When that happens, take it up again the next morning without
fuss, and go on.

Much of a rule can be kept together. The Angelus before the evening meal, the
O antiphon read aloud in the last week, the Sunday Gospel read on Saturday
night, a gift to the poor chosen together: these suit a household, and children
can take part in all of them. A fast that suits the adults should be explained
to the children and made smaller for them, not imposed.

Keep the Sundays as feasts. The Sundays of Advent are the Lord's days in a
season of waiting, and the rule may be relaxed on them. So may 8 December, the
feast of the Immaculate Conception: the Church keeps it with joy, and so should
the household.
